%!TEX root = ../../springer-mono-format/book.tex
%!TEX root = ../../thesis-trier.tex

\IN{permutation-reset automaton}
\IN{permutation automaton}
\IN{reset automaton}
Let $\semiautomat{A}{Q}{\Sigma}{\delta}$.
We call $\mathcal A$ a
\emph{permutation-reset automaton}  
if 	$\delta(Q, a) \ne Q$ implies $|\delta(Q, a)| = 1$
for every $a \in \Sigma$,
a \emph{permutation automaton} if every letter is a permutational letter,
and a \emph{reset automaton} if every letter synchronizes the automaton.
The automaton~$\mathcal A$ is a permutation automaton if and only if~$\TS{\mathcal A}$ is a permutation group on~$Q$.
Permutation automata are SI-automata
without idempotent letters. Clearly, if $\mathcal A$ is a permutation-reset automaton, then
we have
\begin{equation} \label{prelim:perm:fact:synA}
	\Syn(\mathcal A) = \Sigma^* \cdot \{ a \in \Sigma : |\delta(Q, a)| = 1 \} \cdot \Sigma^*
\end{equation}


\begin{proposition} \label{prelim:perm-reset-Rees-factor}
Let $\semiautomat{A}{Q}{\Sigma}{\delta}$ and $S \subseteq Q$ be an upper set.
	Then the following are equivalent.
\begin{enumerate}
\item \label{prelim:perm-reset-Rees-factor:1} $\mathcal A / {\sim_S}$ is a permutation-reset automaton 

\item \label{prelim:perm-reset-Rees-factor:2} $
\delta(Q \setminus S, a) \ne Q \setminus S
$
implies
$ 
\delta(Q, a) \subseteq S
$
for each $a \in \Sigma$.

\item \label{prelim:perm-reset-Rees-factor:3} $Q \setminus S$ is a block.
	
\item \label{prelim:perm-reset-Rees-factor:4} $\mathcal A_{|Q\setminus S, \Gamma}$ for $\Gamma = \{ a \in \Sigma \mid \applytrans{\delta}{Q \setminus S}{a} \subseteq Q \setminus S\}$ is a permutation automaton and $a \in \Sigma\setminus \Gamma$ implies $\applytrans{\delta}{Q \setminus S}{a} \subseteq S$.
\end{enumerate}
\end{proposition}
\begin{proof}
$\ref{prelim:perm-reset-Rees-factor:1} \Rightarrow \ref{prelim:perm-reset-Rees-factor:2}$:
By \cref{prelim:rem:Rees}, the state set~$Q / {\sim_S}$ is permuted by a letter~$a \in \Sigma$
if and only if $\delta(Q \setminus S, a) = Q \setminus S$
and $\delta_{\sim_S}(S, a) = S$.
So, $\delta(Q \setminus S, a) \ne Q \setminus S$
implies $| \delta_{\sim_S}(Q / {\sim_S}, a)| = 1$.
Since~$S$ is a trap state in~$\mathcal A / {\sim_S}$,
this yields $\delta_{\sim_S}(Q / {\sim_S}, a) = \{S\}$,
which is equivalent to $\delta(Q, a) \subseteq S$.

$\ref{prelim:perm-reset-Rees-factor:2} \Rightarrow \ref{prelim:perm-reset-Rees-factor:3}$:
Since~$S$ is an upper set, $\delta(Q, a) \subseteq S$
is equivalent to $\delta(Q \setminus S, a) \subseteq S$.
So, by assumption,
we have $\delta(Q \setminus S, a) = Q \setminus S$
or $\delta(Q \setminus S, a) \subseteq S$,
and the latter condition is equivalent to $\delta(Q \setminus S, a) \cap (Q \setminus S) = \emptyset$.
Thus, $Q \setminus S$ is a block.

$\ref{prelim:perm-reset-Rees-factor:3} \Rightarrow \ref{prelim:perm-reset-Rees-factor:4}$: 
Since~$Q \setminus S$ is a block,
$\applytrans{\delta}{Q \setminus S}{a} \subseteq Q \setminus S$
implies $\applytrans{\delta}{Q \setminus S}{a} = Q \setminus S$, and so $\mathcal A_{|Q\setminus S, \Gamma}$
is a permutation automaton. For $a \in \Sigma \setminus \Gamma$, we have $\delta(Q \setminus S, a) \cap (Q \setminus S) = \emptyset$,
which implies $\delta(Q \setminus S, a) \subseteq S$.

$\ref{prelim:perm-reset-Rees-factor:4} \Rightarrow \ref{prelim:perm-reset-Rees-factor:1}$: 
The assumption yields that $\delta(Q \setminus S, a) = Q \setminus S$
or $\delta(Q \setminus S, a) \subseteq S$ for each $a \in \Sigma$.
Since~$S$ is an upper set, the latter condition
is equivalent to $\delta(Q, a) \subseteq S$.
From \cref{prelim:rem:Rees} we get that $\mathcal A / {\sim_S}$
is a permutation-reset automaton.
\end{proof}

\begin{remark}
	By \cref{prelim:perm-reset-Rees-factor},
	a strongly connected component $S \subseteq Q$ in $\semiautomat{A}{Q}{\Sigma}{\delta}$
	is a block if and only if $( \mathcal A_{|T} ) / {\sim_{T \setminus S}}$
	for $T = \upperset{S}$ is a permutation-reset automaton.	
\end{remark}

By \cref{prelim:perm-reset-Rees-factor} in a permutation-reset automaton
with a trap state $t$, by considering $S = \{t\}$, every non-permutational
letter maps everything to the trap state.

The equivalences from \cref{prelim:perm-reset-Rees-factor} 
fail for $S = \emptyset$ in case of permutation-reset automata $\semiautomat{A}{Q}{\Sigma}{\delta}$
with $|Q| \ge 2$ and a letter $a \in \Sigma$ such that $|\delta(Q, a)| = 1$.

\begin{proposition} \label{prelim:subclasses:perm:reach} 
	Let $\semiautomat{A}{Q}{\Sigma}{\delta}$ be a permutation automaton.
	Then $\delta(q, u) \eqreach q$ for all $q \in Q$
	and $u \in \Sigma^*$.		
\end{proposition}
\begin{proof}
	This follows from \cref{prelim:autsyncw:lem:perm-each-letter-sink} 
	since the whole state set is permuted by each letter.
\end{proof}

\begin{corollary} \label{prelim:subclasses:perm:reach:cor} 
	In a permutation automaton each strongly connected component is a sink component.
\end{corollary}

By \cref{prelim:subclasses:perm:reach}, the strongly connected components in a permutation automaton are precisely 
the orbits of the group $\TS{\mathcal A}$.
Furthermore, for a permutation automaton 
the following are equivalent: strong connectivity, transitivity of the group of units of the transformation monoid,
univalent connectedness, and the property of being cyclic.


Next, using \THM~\ref{prelim:thm:char_minaut}, we prove the following
minimality condition for permutation automata (the special cases $|F| = 1$ or $|F| = |Q| - 1$
were mentioned in~\cite[Cor.~3.8]{DBLP:journals/mst/Davies18}).

\begin{proposition}
	\label{prelim:subclasses:prop:perm:min}
	Let $\automat{A}{Q}{\Sigma}{\delta}{q_0}{F}$ be a permutation automaton with $1 \le |F| < |Q|$
	such that $\gcd(|F|, |Q|) = 1$.
	Then $\mathcal A$ is minimal %for $L(\mathcal A)$
	if and only if $\mathcal A$ is accessible.
\end{proposition}
\begin{proof}
	A non-accessible automaton is not minimal by \cref{prelim:thm:char_minaut}.
	Next, assume $\mathcal A$ is accessible.
	Set $G = \TS{\mathcal A}$. 
	We show that~$\mathcal A$ admits no non-trivial $G$-congruence
	that saturates $F$, which gives minimality of~$\mathcal A$ 
	by \crefitem{prelim:thm:char_minaut}{thm:minaut:item:congruence}.
	Assume~$\sim$ is some $G$-congruence that saturates~$F$.
	Then $F$ is a union of blocks. Let $\Delta \subseteq F$ be a block.
	As $\mathcal A$ is accessible,~$G$ 
	is transitive. Then, by \cref{prelim:ex:trans-block}, $F$
	and $Q$ are unions of equal-sized sets of the form $g(\Delta)$ for $g \in G$. 
	Thus $|\Delta|$ divides~$|F|$ and~$|Q|$, and so it divides $\gcd(|F|, |Q|)$.
	By coprimality, we find $|\Delta| = 1$. So~$\Delta$ is a trivial block
	and the only $G$-congruence saturating $F$ is the equality relation. 
\end{proof}

\begin{corollary}
	Let $\automat{A}{Q}{\Sigma}{\delta}{q_0}{F}$ be a permutation automaton with $1 \le |F| < |Q|$
	such that $\gcd(|F|, |Q|) = 1$.
	If $\TS{\mathcal A}$ contains a transitive permutation group $G$,
	then $\mathcal A$ is minimal.
\end{corollary}
\begin{proof}
    Let $\Gamma = \{ x \in \Sigma \mid \text{$x$ is a permutational letter} \}$.
    Since $G \subseteq \TS{\mathcal A_{|\Gamma}}$,
    $\TS{\mathcal A_{|\Gamma}}$ is a transitive permutation group,
    which yields that $\mathcal A_{|\Gamma}$
    is an accessible permutation automaton.
    So, by \cref{prelim:subclasses:prop:perm:min}, 
    $\mathcal A_{|\Gamma}$ is minimal.
    Since accessibility and distinguishability
    of pairs of states extend
    from~$\mathcal A_{|\Gamma}$ to $\mathcal A$,
    also~$\mathcal A$ is minimal by \crefitem{prelim:thm:char_minaut}{thm:minaut:item:statesdist}.
\end{proof}